\chapter{Before we start}

\section{What raylib actually is}

raylib is a C library that hands you a window, a keyboard, a mouse, a sound card
and a way to put coloured pixels on screen. That is it. It is not an engine.
There is no editor, no scene graph, no entity system, no asset pipeline, no
project format. You write a \lstinline|main()|, you call some functions, you get
a game.

That is the reason it feels strange at first if you are used to tools that hand
you a scene and ask you to fill it in. In raylib \emph{you} own the loop. Nothing
happens that you did not write a line for.

The whole library is about 500 functions, and they are named so plainly that you
can usually guess them: \lstinline|DrawCircle|, \lstinline|IsKeyDown|,
\lstinline|LoadTexture|, \lstinline|CheckCollisionRecs|. There is no official
reference manual, on purpose --- the header file \file{raylib.h} \emph{is} the
documentation, one line of comment per function, and the examples are the
tutorial.

\keyidea{If you remember one thing from this manual, remember this: the
\file{raylib.h} header is the documentation. Open it. Search it. Every function
you will ever use is in there with a one-line comment explaining it.}

\section{What you need to know already}

C. Not advanced C --- variables, \lstinline|if|, \lstinline|while|, functions,
arrays and structs are enough. Every piece of code in this manual stays inside
that. You do not need pointers-to-pointers, you do not need
\lstinline|malloc|, and you certainly do not need C++.

You do \emph{not} need any 3D background. Chapters 9 to 14 assume you have never
seen a 3D coordinate system in your life, and build it up from a single arrow on
a page.

\section{The files that come with this manual}

This manual is not a thing to read in an armchair. Every chapter has a matching
program you can run, change and break. They live in two folders next to it:

\begin{center}
\begin{tabular}{@{}lll@{}}
\toprule
\textbf{Folder} & \textbf{Part} & \textbf{What is in it} \\
\midrule
\file{course2d/} & I & the fundamentals: window, loop, shapes, input, collisions, camera \\
\file{course3d/} & I & the 3D half: axes, cameras, movement, collisions, the mini Doom \\
\file{cheatsheet/} & II & one example per module of the official raylib cheatsheet \\
\bottomrule
\end{tabular}
\end{center}

Each folder has a \file{build.sh} that compiles and runs one of them:

\begin{lstlisting}[style=plain]
./course2d/build.sh 1        # builds and runs course2d/lesson01_window.c
./course3d/build.sh 7        # builds and runs the mini Doom
./cheatsheet/build.sh 9      # builds and runs the audio module example
\end{lstlisting}

Part I is a path: read it in order and you end up with a game. Part II is a
reference: each chapter stands alone, so read the one you need when you need it.

The script needs the raylib library to have been built once. If you have never
run the project, do that first from the root folder:

\begin{lstlisting}[style=plain]
make                         # builds bin/Debug/libraylib.a and the sample app
\end{lstlisting}

\section{What a build actually needs}

When people get stuck on day one, it is almost never the code --- it is the
compiler command. There is no magic in it, so here it is in full, the one that
\file{build.sh} runs for you:

\begin{lstlisting}[style=plain]
cc  mygame.c                                   # your source file
    -o mygame                                  # the program to produce
    -I build/external/raylib-master/src        # where raylib.h lives
    bin/Debug/libraylib.a                      # the library itself
    -lpthread -lm -ldl -lrt -lX11              # what raylib needs on Linux
\end{lstlisting}

Three ingredients, always: your \file{.c} file, the \emph{header} so the compiler
knows the function signatures, and the \emph{library} so the linker can find the
actual code. If you get \lstinline|raylib.h: No such file| you got the
\lstinline|-I| wrong. If you get \lstinline|undefined reference to InitWindow|
you got the library path wrong. Those two errors mean two very different things
and mixing them up costs beginners hours.

\gotcha{The order matters on Linux: the library goes \emph{after} your
\file{.c} file on the command line, never before. GNU \lstinline|ld| resolves
symbols left to right, and a library listed first has nothing left to resolve.}

\section{How to read this manual}

Read a chapter, then run its lesson, then go back and change three numbers in the
lesson until something breaks. That third step is the one that actually teaches
you. Nothing here can be damaged: the worst that happens is a black window, and
you press Escape.

The boxes mean:

\gotcha{A mistake almost everyone makes here. Reading these will save you an
afternoon each.}

\keyidea{The one sentence from the section that is worth carrying with you.}

\tryit{The lesson that shows this chapter running.}{./course2d/build.sh 1}
